\documentclass[10pt, aspectratio=169]{beamer}
\usepackage[utf8]{inputenc}
\usepackage[portuguese]{babel}
\usepackage{amsmath, amssymb, amsfonts}
\usepackage{graphicx}
\usepackage{booktabs}
\usepackage{tikz}
\usetikzlibrary{shapes, arrows.meta, positioning, calc}
\usepackage{listings}
\usepackage{xcolor}

% Color Palette - Professional Data Science & Interactive Class Theme
\definecolor{navyblue}{RGB}{30, 41, 59}       % #1E293B Primary dark
\definecolor{skycyan}{RGB}{14, 165, 233}      % #0EA5E9 Secondary blue/cyan
\definecolor{ambergold}{RGB}{245, 158, 11}    % #F59E0B Accent amber
\definecolor{darkgray}{RGB}{51, 65, 85}       % Dark text / subheaders
\definecolor{lightgray}{RGB}{241, 245, 249}   % Card / block background
\definecolor{codebg}{RGB}{248, 249, 250}      % Code background
\definecolor{codegreen}{RGB}{34, 197, 94}     % Code comments/strings
\definecolor{crimsonred}{RGB}{225, 29, 72}     % Highlight red
\definecolor{livepurple}{RGB}{147, 51, 234}    % Purple for Live Coding

% Beamer Theme Settings
\usetheme{Madrid}
\usecolortheme{whale}

\setbeamercolor{structure}{fg=navyblue}
\setbeamercolor{palette primary}{bg=navyblue,fg=white}
\setbeamercolor{palette secondary}{bg=skycyan,fg=white}
\setbeamercolor{palette tertiary}{bg=darkgray,fg=white}
\setbeamercolor{titlelike}{bg=navyblue,fg=white}

\setbeamercolor{block title}{bg=navyblue,fg=white}
\setbeamercolor{block body}{bg=lightgray,fg=black}

\setbeamercolor{block title alert}{bg=crimsonred,fg=white}
\setbeamercolor{block body alert}{bg=lightgray,fg=black}

\setbeamercolor{block title example}{bg=skycyan,fg=white}
\setbeamercolor{block body example}{bg=lightgray,fg=black}

% Custom Block for Live Coding
\newenvironment{liveblock}[1]{%
    \setbeamercolor{block title}{bg=livepurple,fg=white}%
    \setbeamercolor{block body}{bg=lightgray,fg=black}%
    \begin{block}{#1}%
}{%
    \end{block}%
}

\setbeamertemplate{navigation symbols}{}
\setbeamertemplate{footline}[frame number]

% Python Listing Setup
\lstset{
    language=Python,
    backgroundcolor=\color{codebg},
    basicstyle=\ttfamily\scriptsize\color{navyblue},
    keywordstyle=\bfseries\color{skycyan},
    stringstyle=\color{codegreen},
    commentstyle=\itshape\color{gray},
    numbers=none,
    breaklines=true,
    frame=single,
    rulecolor=\color{lightgray}
}

% Title Slide Metadata
\title[Ciência de Dados: Probabilidade e Média]{Ciência de Dados -- Aula 08 (Parte 1)}
\subtitle{Revisão Aprofundada de Probabilidade, Amostragem e Propriedades da Média\\{\small Guia Explicativo Detalhado de Equações e Estudos de Caso}}
\author[ADS]{Curso de Análise e Desenvolvimento de Sistemas (ADS)}
\institute[ADS]{Disciplina: Ciência de Dados}
\date{\today}

\begin{document}

% -----------------------------------------------------------------------------
% Slide 1: Title
% -----------------------------------------------------------------------------
\begin{frame}
    \titlepage
\end{frame}

% -----------------------------------------------------------------------------
% Slide 2: Agenda
% -----------------------------------------------------------------------------
\begin{frame}{Agenda Explicativa da Parte 1}
    \tableofcontents
\end{frame}

% =============================================================================
\section{1. Revisão Detalhada de Probabilidade para ADS}
% =============================================================================

\begin{frame}{1.1 Espaço Amostral, Experimentos e Eventos}
    \begin{block}{Definições Fundamentais em Ciência de Dados}
        \begin{itemize}
            \item \textbf{Experimento Aleatório:} Qualquer ação ou processo observável cujo resultado final não pode ser determinado com certeza absoluta antes da execução (ex: envio de requisição web, sorteio de usuário).
            \item \textbf{Espaço Amostral ($\Omega$):} O conjunto de \textbf{todos} os resultados possíveis do experimento.
            \item \textbf{Evento ($A$):} Qualquer subconjunto do espaço amostral ($A \subseteq \Omega$).
        \end{itemize}
    \end{block}

    \vspace{0.2cm}

    \begin{exampleblock}{Fórmula da Probabilidade Clássica (Resultados Equiprováveis)}
        \[ P(A) = \frac{|A|}{|\Omega|} = \frac{\text{Número de resultados favoráveis ao evento } A}{\text{Número total de resultados possíveis no espaço amostral } \Omega} \]
        \textbf{Significado das Variáveis e Símbolos:}
        \begin{itemize}
            \item $P(A)$: Probabilidade de ocorrência do evento $A$ (número real entre $0$ e $1$).
            \item $|A|$: Cardinalidade do conjunto $A$ (quantidade de elementos no evento $A$).
            \item $|\Omega|$: Cardinalidade do espaço amostral $\Omega$ (quantidade total de resultados possíveis).
        \end{itemize}
    \end{exampleblock}
\end{frame}

\begin{frame}{1.2 Regra da Adição e União de Eventos: $P(A \cup B)$}
    \begin{block}{Eventos Mutuamente Exclusivos}
        Dois eventos $A$ e $B$ são \textbf{mutuamente exclusivos} (ou disjuntos) se nunca puderem ocorrer simultaneamente no mesmo experimento ($A \cap B = \emptyset$). 
        \[ P(A \cup B) = P(A) + P(B) \]
    \end{block}
    \begin{alertblock}{Regra Geral da Adição (Eventos Não-Exclusivos)}
        Quando existe sobreposição entre os eventos ($A \cap B \neq \emptyset$):
        \vspace{-0.2cm}
        \[ P(A \cup B) = P(A) + P(B) - P(A \cap B) \]
        \textbf{Explicação das Variáveis:}
        \begin{itemize}
            \item $P(A \cup B)$: Probabilidade de ocorrer o evento $A$ \textbf{OU} o evento $B$ (União de conjuntos).
            \item $P(A)$: Probabilidade individual de ocorrência do evento $A$.
            \item $P(B)$: Probabilidade individual de ocorrência do evento $B$.
            \item $P(A \cap B)$: Probabilidade de ocorrerem \textbf{AMBOS} os eventos ao mesmo tempo (Interseção).
            \item \textbf{Por que subtrair $P(A \cap B)$?} Na soma $P(A) + P(B)$, a região de interseção é contada \textbf{duas vezes}. Subtraímos uma vez para corrigir a contagem dupla!
        \end{itemize}
    \end{alertblock}
\end{frame}

\begin{frame}{1.3 Probabilidade Condicional e Regra da Multiplicação}
	\vspace{-0.2cm}
    \begin{block}{Regra Geral da Multiplicação}
        Para quaisquer dois eventos $A$ e $B$, a probabilidade de ambos ocorrerem é:
        \vspace{-0.2cm}
        \[ P(A \cap B) = P(A) \cdot P(B|A) \]
        \textbf{Explicação Detalhada das Variáveis:}
        \begin{itemize}
            \item $P(A \cap B)$: Probabilidade de ocorrer $A$ \textbf{E} $B$ em sequência ou simultaneamente.
            \item $P(A)$: Probabilidade do primeiro evento $A$ ocorrer.
            \item $P(B|A)$: Probabilidade condicional de $B$ ocorrer \textbf{dado que $A$ já ocorreu}.
        \end{itemize}
    \end{block}

    \begin{exampleblock}{Estudo de Caso Explicado: Sorteio Sem Reposição}
        \textbf{Problema:} Em um grupo de 3 alunos $\{X, Y, Z\}$, qual a probabilidade de sortear $Y$ no 1º turno e $X$ no 2º turno sem reposição?
        \begin{itemize}
            \item \textbf{Passo 1 (Evento $A$):} Probabilidade de tirar $Y$ primeiro $\implies P(A) = \frac{1}{3}$.
            \item \textbf{Passo 2 (Evento $B|A$):} Com $Y$ fora, restam 2 alunos $\{X, Z\}$. Probabilidade de tirar $X$ em seguida $\implies P(B|A) = \frac{1}{2}$.
            \item \textbf{Solução e Resposta Final:} $P(A \cap B) = \frac{1}{3} \cdot \frac{1}{2} = \frac{1}{6} \approx 0{,}1667 \text{ (ou } 16{,}67\%\text{)}$.
        \end{itemize}
    \end{exampleblock}
\end{frame}

\begin{frame}{1.4 Independência de Eventos e a Crise Financeira de 2008}
    \begin{block}{Definição Formal de Independência}
        Dois eventos $A$ e $B$ são \textbf{independentes} se a ocorrência de $A$ não altera a probabilidade de $B$ ocorrer ($P(B|A) = P(B)$).
        \[ P(A \cap B) = P(A) \cdot P(B) \quad \text{(Válido APENAS para eventos independentes)} \]
    \end{block}

    \vspace{0.2cm}

    \begin{alertblock}{Estudo de Caso Explicado: A Falha da Suposição de Independência em 2008}
        \textbf{O Problema Real:} Durante a crise dos títulos imobiliários de 2008, modelos financeiros assumiram que a inadimplência de hipotecas em diferentes cidades eram eventos \textbf{independentes}.
        \begin{itemize}
            \item \textbf{Cálculo Errado do Mercado:} Se $P(\text{Inadimplência em Cidade A}) = 5\%$ e em Cidade B = $5\%$, calcularam a probabilidade de falha dupla como:
                  \[ P(A \cap B) = 0{,}05 \times 0{,}05 = 0{,}0025 \quad (0{,}25\% \text{ -- Risco quase zero!}) \]
            \item \textbf{Como Chegamos à Resposta Real:} As economias estavam interligadas (fatores sistêmicos). A inadimplência na Cidade A aumentava a chance de inadimplência na Cidade B ($P(B|A) \gg P(B)$). A falta de independência causou o colapso dos modelos!
        \end{itemize}
    \end{alertblock}
\end{frame}

\begin{frame}[fragile]{[LIVE CODING] Prática 1: Probabilidade em Python}
    \begin{liveblock}{[PRÁTICA EM SALA] Execução no Notebook (Célula 2)}
        Simulando o cálculo empírico de probabilidade condicional e união:
    \end{liveblock}

\begin{lstlisting}
import numpy as np

# Simulando 10.000 lancamentos de 2 moedas
np.random.seed(42)
m1 = np.random.choice(['H', 'T'], size=10000)
m2 = np.random.choice(['H', 'T'], size=10000)

# Evento: Pelo menos uma cara (H) -> P = 3/4 = 75%
p_pelo_menos_uma_cara = np.mean((m1 == 'H') | (m2 == 'H'))
print(f"P(Pelo menos 1 Cara) observada: {p_pelo_menos_uma_cara:.2%}")

# Multiplicacao sem reposicao (1/3 * 1/2 = 1/6)
p_yx = (1/3) * (1/2)
print(f"P(Y 1o e X 2o sem reposicao): {p_yx:.4f} (Teorico: 0.1667)")
\end{lstlisting}
\end{frame}

% =============================================================================
\section{2. Amostragem Aleatória em Python}
% =============================================================================

\begin{frame}{2.1 População vs. Amostra e Tipos de Amostragem}
    \begin{block}{Parâmetro Populacional vs. Estatística Amostral}
        \begin{itemize}
            \item \textbf{População:} O conjunto total e absoluto de todos os elementos estudados ($N$ maiúsculo). Parâmetro populacional é fixo (ex: média populacional $\mu$).
            \item \textbf{Amostra:} Um subconjunto aleatório extraído da população ($n$ minúsculo). Estatística amostral varia a cada amostra (ex: média amostral $\bar{x}$).
        \end{itemize}
    \end{block}

    \vspace{0.2cm}

    \begin{columns}[T]
        \begin{column}{0.48\textwidth}
            \begin{exampleblock}{Amostragem COM Reposição}
                Cada elemento sorteado retorna ao conjunto antes do próximo sorteio.
                \begin{itemize}
                    \item Resultados são \textbf{independentes}.
                    \item Espaço amostral não muda de tamanho.
                \end{itemize}
            \end{exampleblock}
        \end{column}

        \begin{column}{0.48\textwidth}
            \begin{block}{Amostragem SEM Reposição}
                Cada elemento sorteado é removido permanentemente.
                \begin{itemize}
                    \item Resultados são \textbf{independentes} (probabilidades mudam).
                    \item Utilizado em amostragem aleatória simples de bancos de dados.
                \end{itemize}
            \end{block}
        \end{column}
    \end{columns}
\end{frame}

\begin{frame}[fragile]{[LIVE CODING] Prática 2: Amostragem com NumPy e Pandas}
    \begin{liveblock}{[PRÁTICA EM SALA] Amostragem de DataFrames}
        Executando amostragem aleatória simples em conjuntos de dados:
    \end{liveblock}

\begin{lstlisting}
import pandas as pd

populacao = np.array(['Aluno_A', 'Aluno_B', 'Aluno_C', 'Aluno_D', 'Aluno_E'])

# Sorteio COM reposicao
print("COM reposicao:", np.random.choice(populacao, size=3, replace=True))

# Sorteio SEM reposicao
print("SEM reposicao:", np.random.choice(populacao, size=3, replace=False))

# Amostragem direta de um DataFrame no Pandas
df_users = pd.DataFrame({'id': range(100), 'ativo': np.random.choice([0,1], 100)})
amostra = df_users.sample(n=5, random_state=42)
print("Amostra DataFrame:\n", amostra)
\end{lstlisting}
\end{frame}

% =============================================================================
\section{3. A Média e suas Propriedades Fundamentais}
% =============================================================================

\begin{frame}{3.1 Média Aritmética como Operação de Suavização (*Smoother*)}
	\vspace{-0.2cm}
    \begin{block}{Fórmula da Média Aritmética Simples}
        \[ \bar{x} = \frac{1}{N} \sum_{i=1}^N x_i = \frac{x_1 + x_2 + \dots + x_N}{N} \]
        \textbf{Significado Explicado das Variáveis:}
        \begin{itemize}
            \item $\bar{x}$ (lê-se "x-barra"): A média aritmética do conjunto de dados.
            \item $N$: Número total de elementos/observações na coleção.
            \item $x_i$: O valor individual do $i$-ésimo elemento da coleção.
            \item $\sum_{i=1}^N$: Símbolo de somatório do primeiro elemento ($i=1$) até o último ($i=N$).
        \end{itemize}
    \end{block}
    \begin{exampleblock}{Conceito do Fundo Comum (Equalizador)}
        Para a coleção $\{2, 3, 3, 9\}$:
        \vspace{-0.4cm}
        \[ \text{Total} = 2 + 3 + 3 + 9 = 17 \implies \bar{x} = \frac{17}{4} = 4{,}25 \]
        Se juntarmos todos os \$17 em um fundo comum e redistribuirmos igualmente, cada indivíduo recebe \$4,25. A média suaviza as disparidades individuais.
    \end{exampleblock}
\end{frame}

\begin{frame}{3.2 Proporções são Médias de Variáveis Binárias}
    \begin{block}{Matemática das Variáveis Indicadoras ($0$ e $1$)}
        Em sistemas de informação, dados binários ou booleanos são representados por $x_i \in \{0, 1\}$.
        \[ \text{Proporção de 1s} = \bar{x} = \frac{\sum_{i=1}^N x_i}{N} = \frac{\text{Número total de 1s (Sucessos)}}{\text{Número total de elementos } N} \]
    \end{block}

    \vspace{0.2cm}

    \begin{exampleblock}{Aplicações Práticas explicadas para ADS}
        \begin{itemize}
            \item \textbf{Taxa de Clique (CTR):} Se um usuário clicou ($1$) ou não ($0$). A média do vetor binário é exatamente o CTR.
            \item \textbf{Acurácia de Modelo de Machine Learning:} Comparação \texttt{(y\_real == y\_pred)} gera um vetor de \texttt{True} ($1$) e \texttt{False} ($0$). A média resulta na acurácia percentual!
        \end{itemize}
    \end{exampleblock}
\end{frame}

\begin{frame}{3.3 O Centro de Gravidade do Histograma}
    \vspace{-0.2cm}
    \begin{block}{Cálculo da Média por Ponderação de Frequências Relativas}
        Agrupando os valores distintos $x_j$ ponderados pelas suas proporções $p_j$:
        \vspace{-0.5cm}
        \[ \bar{x} = \sum_{j=1}^k x_j \cdot p_j \]
        \textbf{Significado Explicado das Variáveis:}
        \begin{itemize}
            \item $x_j$: Cada valor \textbf{distinto} presente no conjunto de dados.
            \item $p_j$: A \textbf{frequência relativa (proporção)} com que o valor $x_j$ aparece na coleção ($p_j = \text{frequência}/N$).
            \item $k$: Quantidade de valores distintos na coleção.
        \end{itemize}
    \end{block}

    \begin{exampleblock}{Demonstração Passo a Passo para $\{2, 3, 3, 9\}$}
        \begin{itemize}
            \item Valor $2$: aparece 1 vez em 4 $\implies p_1 = 1/4 = 0{,}25$.
            \item Valor $3$: aparece 2 vezes em 4 $\implies p_2 = 2/4 = 0{,}50$.
            \item Valor $9$: aparece 1 vez em 4 $\implies p_3 = 1/4 = 0{,}25$.
            \item \textbf{Cálculo Ponderado:} $\bar{x} = (2 \cdot 0{,}25) + (3 \cdot 0{,}50) + (9 \cdot 0{,}25) = 0{,}5 + 1{,}5 + 2{,}25 = 4{,}25$.
        \end{itemize}
    \end{exampleblock}
\end{frame}

% =============================================================================
\section{4. Média vs. Mediana e Assimetria (Skewness)}
% =============================================================================

\begin{frame}{4.1 O Mito do "Abaixo da Média"}
    \begin{alertblock}{Desmistificando a Relação Média vs. Mediana}
        \textbf{Pergunta:} "Se tirar nota 4 em uma prova onde a média foi 4.25, estou na metade inferior da turma?"
        \begin{itemize}
            \item \textbf{Resposta:} \textbf{Não necessariamente!}
            \item \textbf{Explicação Detalhada:} Na distribuição de notas $\{2, 3, 3, 9\}$, a média é $4{,}25$ e a mediana é $3{,}00$. Os elementos $2, 3, 3$ estão todos abaixo da média. Portanto, \textbf{$75\%$ da turma está abaixo da média!}
            \item \textbf{Por que isso ocorre?} A média é o ponto de equilíbrio de massa, não o ponto de corte de 50\% dos dados (que é a mediana).
        \end{itemize}
    \end{alertblock}
\end{frame}

\begin{frame}{4.2 Estudo de Caso Explicado: Salários de San Francisco (SF 2015)}
    \begin{block}{O Problema Real de Negócio e Dados}
        Análise da folha de pagamento de San Francisco em 2015 ($N > 30.000$ servidores).
        \begin{itemize}
            \item \textbf{Mediana Observada:} $\approx \$99.359$
            \item \textbf{Média Observada:} $\approx \$111.943$
            \item \textbf{Diferença:} A média é R\$ 12.584 \textbf{maior} do que a mediana!
        \end{itemize}
    \end{block}

    \vspace{0.2cm}

    \begin{exampleblock}{Como Chegamos à Explicação e Solução}
        \begin{itemize}
            \item \textbf{Causa:} O histograma tem assimetria à direita (*Right Skewed*), com uma cauda longa formada por médicos e altos executivos com salários de até \$600.000.
            \item \textbf{Decisão de ADS:} Para relatar o salário do funcionário público "típico", deve-se usar a \textbf{Mediana} (resistente a outliers). Para planejar o orçamento total da prefeitura, deve-se usar a \textbf{Média}($\text{Orçamento Total} = N \times \text{Média}$).
        \end{itemize}
    \end{exampleblock}
\end{frame}

\begin{frame}{Resumo da Aula 1}
    \begin{block}{Síntese da Aula 1}
        \begin{enumerate}
            \item \textbf{Probabilidade:} Espaço amostral $\Omega$, união $P(A \cup B)$ e produto condicional $P(A \cap B) = P(A)P(B|A)$.
            \item \textbf{Amostragem:} Com reposição (independente) vs. sem reposição (dependente).
            \item \textbf{Média:} Ponto de equilíbrio $\bar{x} = \sum x_j p_j$ e equivalência com proporções para dados binários.
            \item \textbf{Assimetria:} A média desliza na direção da cauda longa; a mediana reflete o ponto central de 50\%.
        \end{enumerate}
    \end{block}
\end{frame}

\end{document}
